% =========================================================================
% Chapter 7: Time-Varying Carbon-Conditional Hazard — A State-Space Approach
% Sake Saakstra, MSc EOR thesis, Vrije Universiteit Amsterdam
% Supervisor: prof. dr. S.J. Koopman
% =========================================================================
\documentclass[11pt,a4paper]{article}

\usepackage[a4paper,margin=2.5cm]{geometry}
\usepackage{amsmath,amssymb,amsthm,bm}
\usepackage{graphicx}
\usepackage{booktabs}
\usepackage{natbib}
\usepackage{hyperref}
\usepackage{xcolor}
\usepackage{caption}
\hypersetup{colorlinks=true,linkcolor=blue!50!black,citecolor=blue!50!black,urlcolor=blue!50!black}

\title{\textbf{Chapter 7}\\[0.3em] Time-Varying Carbon-Conditional Hazard:\\
A State-Space Approach to Hydrogen Project Cancellation}
\author{Sake Saakstra \\ \small Supervisor: prof. dr. S.J. Koopman}
\date{Draft, \today}

\newcommand{\bint}{\beta_{\mathrm{int}}}
\newcommand{\bblue}{\beta_{\mathrm{blue}}}
\newcommand{\beua}{\beta_{\mathrm{eua}}}
\newcommand{\E}{\mathbb{E}}

\begin{document}
\maketitle

% =========================================================================
\section{Introduction and Motivation}
\label{sec:ch7-intro}

The static carbon-conditional hazard model developed in Chapter 6 estimates a single interaction coefficient $\bint$ governing the differential cancellation risk of blue versus green hydrogen projects as a function of the EU Emissions Allowance (EUA) price. The maximum a posteriori estimate of $\bint = -1.43$ with 95\% credible interval $[-2.59, -0.37]$ implies a structurally negative dependence: at higher EUA prices, the cancellation hazard of blue projects falls relative to the green (PEM electrolyser) reference, while at lower EUA prices the gap widens substantially.

This static specification, while informative, embeds a strong implicit assumption: the carbon-conditional sensitivity is constant throughout the 2010--2026 observation window. Several economic developments motivate questioning this assumption. The EUA market itself underwent dramatic transitions, from the surplus-driven era of 2010--2017 (where allowances traded below \euro 10) through the Market Stability Reserve activation in 2019, the supply-tightening induced by the 2022--2023 energy crisis (peaks above \euro 100), and the announcement of CBAM and the ETS2 framework in 2022--2023. Concurrently, the hydrogen project landscape shifted: the share of announced electrolyser capacity rose from below 5\% of global hydrogen production projects in 2018 to over 60\% by 2023, while major blue hydrogen cancellations clustered in 2023--2024 (79\% of all events in our sample).

If the underlying decision-making behaviour of project sponsors adapted to these shifts---for example, by attaching different weights to carbon prices once Phase IV of the EU ETS clarified the free-allocation phase-out trajectory---then the marginal effect of EUA prices on blue-versus-green relative survival should itself be time-varying. Failing to model this would mask economically meaningful regime dependence and may bias inferences when extrapolated to future scenarios. The methodological tools for testing this exist in the state-space literature \citep{durbin2012time, koopman2000time}, but to our knowledge have not been applied to project-level survival in the energy transition.

We therefore develop a sequence of three nested specifications, each relaxing assumptions about the temporal structure of $\bint$:
\begin{enumerate}
    \item[(M1)] \textbf{Static}: $\bint(t) = \bint$ for all $t$.
    \item[(M2)] \textbf{Parameter-driven TVP}: $\bint(t)$ follows a Gaussian random walk over economically motivated time blocks.
    \item[(M3)] \textbf{Observation-driven TVP (GAS)}: $\bint(t)$ updates according to the score of the likelihood, following \citet{creal2008general, creal2013generalized}.
\end{enumerate}
The progression from M1 to M3 mirrors the parameter-driven vs.\ observation-driven distinction developed in \citet{koopman2016predicting}, who show that the two paradigms offer complementary strengths: parameter-driven models can capture genuine stochastic variation but require informative priors when data are sparse, while observation-driven models regularise the trajectory through structured persistence parameters at the cost of less direct probabilistic interpretation.

This chapter contributes to the literature on three fronts. First, methodologically, we provide a direct comparison of the three specifications on a common dataset and discuss how their findings diverge or converge. Second, substantively, we provide evidence on the temporal stability of the carbon-conditional cancellation mechanism in hydrogen project investments, complementing recent work on the financial economics of the energy transition \citep{bolton2021carbon, odenweller2022probabilistic}. Third, the chapter develops a Bayesian inference framework via Hamiltonian Monte Carlo for both parameter-driven and observation-driven TVP models in discrete-time hazard analysis, which is to our knowledge novel in this application context.

The remainder of the chapter proceeds as follows. Section \ref{sec:ch7-spec} formalises the three specifications. Section \ref{sec:ch7-data} describes the data aggregation choices required for comparable inference across the three models. Section \ref{sec:ch7-results} presents the main findings. Section \ref{sec:ch7-robust} reports robustness checks across alternative block boundaries and GAS scaling choices. Section \ref{sec:ch7-discussion} discusses the substantive and methodological implications. Section \ref{sec:ch7-conclusion} concludes.

% =========================================================================
\section{Specifications}
\label{sec:ch7-spec}

Let $i = 1, \dots, n$ index hydrogen projects, with binary technology indicator $\mathrm{Blue}_i \in \{0,1\}$ taking value one for blue (CCS-based) projects and zero for green (PEM electrolyser) projects. Let $t = t_0, \dots, T$ index calendar years over the observation window 2010--2026. Each project enters at announcement year $t_i^{\mathrm{start}}$ and is at risk until the realisation of a cancellation event or right-censoring. Let $y_{it} = 1$ if project $i$ is cancelled in year $t$, $y_{it} = 0$ otherwise (and undefined if project $i$ is not at risk in year $t$). The standardised EUA price in year $t$ is denoted $z_t$ (with $z_t = (\mathrm{EUA}_t - \bar{\mathrm{EUA}})/\hat\sigma_{\mathrm{EUA}}$).

\subsection{M1: Static baseline}

The static specification treats all parameters as time-invariant. With $\eta_{it}$ the logit-linear predictor and $h_{it} = \mathrm{logit}^{-1}(\eta_{it})$ the discrete-time hazard:
\begin{align}
\eta_{it} &= \alpha + \bblue \cdot \mathrm{Blue}_i + \beua \cdot z_t + \bint \cdot \mathrm{Blue}_i \cdot z_t, \\
y_{it} &\sim \mathrm{Bernoulli}(h_{it}).
\end{align}
At the aggregate level, summing over individuals within (year, technology) cells yields the equivalent Binomial representation:
\begin{align}
n_{t}^{\mathrm{blue,events}} &\sim \mathrm{Binomial}(n_{t}^{\mathrm{blue,atrisk}}, p_{t}^{\mathrm{blue}}), \\
n_{t}^{\mathrm{pem,events}} &\sim \mathrm{Binomial}(n_{t}^{\mathrm{pem,atrisk}}, p_{t}^{\mathrm{pem}}),
\end{align}
with $p_t^{\mathrm{blue}} = \mathrm{logit}^{-1}(\alpha + \bblue + \beua z_t + \bint z_t)$ and $p_t^{\mathrm{pem}} = \mathrm{logit}^{-1}(\alpha + \beua z_t)$. The aggregate representation is the working basis for all three specifications in this chapter, both for computational tractability under HMC and to ensure direct comparability via information criteria.

We place the following weakly informative priors:
\begin{align}
\alpha &\sim \mathcal{N}(-4, 1), &
\bblue &\sim \mathcal{N}(0, 2), \\
\beua &\sim \mathcal{N}(0, 2), &
\bint &\sim \mathcal{N}(0, 2).
\end{align}

\subsection{M2: Parameter-driven TVP via Gaussian random walk over blocks}

The parameter-driven specification allows $\bint$ to vary across $B$ economically motivated time blocks. With $b(t) \in \{0, \dots, B-1\}$ the block index of year $t$:
\begin{align}
\eta_{it} &= \alpha + \bblue \cdot \mathrm{Blue}_i + \beua \cdot z_t + \bint^{b(t)} \cdot \mathrm{Blue}_i \cdot z_t, \\
\bint^{b+1} &= \bint^{b} + \eta_b^{\mathrm{int}}, \qquad \eta_b^{\mathrm{int}} \sim \mathcal{N}(0, \sigma^2_{\mathrm{int}}), \\
\bint^{0} &\sim \mathcal{N}(-1, 2), \\
\sigma_{\mathrm{int}} &\sim \mathrm{HalfNormal}(0.5).
\end{align}
We use $B = 4$ as the baseline specification, with blocks corresponding to (i) pre-energy-crisis (2010--2019), (ii) pandemic and early energy-crisis (2020--2022), (iii) peak cancellations (2023--2024), and (iv) post-peak normalisation (2025--2026). Section \ref{sec:ch7-robust} reports results for $B \in \{3, 4, 5\}$ as a sensitivity check.

This specification corresponds to a discrete-time non-Gaussian state-space model in the sense of \citet{durbin2012time}: the state vector $\bint^{0:B-1}$ follows a Gaussian random walk, while the observation equation is Binomial (and hence non-Gaussian). Bayesian inference via Hamiltonian Monte Carlo allows full posterior characterisation of the latent trajectory, side-stepping the importance sampling routines typically required for likelihood evaluation in such models \citep[Chapter 11]{durbin2012time}.

\subsection{M3: Observation-driven TVP via Generalised Autoregressive Score}

The observation-driven specification follows \citet{creal2008general, creal2013generalized}, in which the time-varying parameter updates as a function of the score of the conditional log-likelihood. For the carbon-conditional coefficient at annual frequency:
\begin{align}
\bint(t+1) &= \omega (1 - \phi) + \phi \, \bint(t) + \alpha_{\mathrm{gas}} \cdot \tilde{s}_t, \\
s_t &= \frac{\partial \log L_t}{\partial \bint(t)} = \big(n_t^{\mathrm{blue,events}} - n_t^{\mathrm{blue,atrisk}} \cdot p_t^{\mathrm{blue}}\big) \cdot z_t, \\
\tilde{s}_t &= \mathcal{I}_t^{-d} \cdot s_t, \\
\mathcal{I}_t &= n_t^{\mathrm{blue,atrisk}} \cdot p_t^{\mathrm{blue}} (1 - p_t^{\mathrm{blue}}) \cdot z_t^2.
\end{align}
The score $s_t$ is the gradient of the conditional log-likelihood at $\bint(t)$; intuitively, it captures how much the model would adjust $\bint(t)$ to better fit the observation in year $t$. The scaling exponent $d$ controls how heavily the score is normalised by the Fisher information at $t$: $d=0$ uses the raw score, $d=1/2$ provides a square-root scaling (our baseline), and $d=1$ applies full information scaling. The persistence parameter $\phi \in [0,1]$ governs the rate of mean reversion toward the long-run mean $\omega$, while $\alpha_{\mathrm{gas}} \geq 0$ scales the response of $\bint(t)$ to news in the score.

Priors for the GAS hyperparameters:
\begin{align}
\omega &\sim \mathcal{N}(-1, 1), &
\phi &\sim \mathrm{Beta}(8, 2), \\
\alpha_{\mathrm{gas}} &\sim \mathrm{HalfNormal}(0.5), &
\bint(0) &\sim \mathcal{N}(-1, 1.5).
\end{align}
The $\mathrm{Beta}(8,2)$ prior on $\phi$ encodes a prior preference for persistent dynamics with prior mean $0.8$, consistent with empirical findings on slowly-evolving structural parameters \citep{creal2013generalized}.

A central methodological feature of the GAS specification is that, given the hyperparameters $(\omega, \phi, \alpha_{\mathrm{gas}}, \bint(0))$ and the observed data, the trajectory $\bint(t)$ is \textit{deterministic}: it is computed by forward iteration of the recursion. The latent trajectory is therefore not a separate random variable to be sampled, but a deterministic function of the four scalar hyperparameters and the data history. This makes the inference problem substantially more parsimonious than parameter-driven alternatives, at the cost of constraining the trajectory to belong to the family generated by the score-driven recursion.

% =========================================================================
\section{Data Aggregation and Implementation}
\label{sec:ch7-data}

For comparability across the three specifications via information criteria, we aggregate the person-year panel of $N = 4{,}162$ observations (covering 714 projects over 2010--2026) to a year-by-technology level. For each year $t$ and technology $\tau \in \{\mathrm{blue}, \mathrm{pem}\}$, we compute $n_t^{\tau,\mathrm{atrisk}}$ (the number of projects at risk in year $t$) and $n_t^{\tau,\mathrm{events}}$ (the number of cancellation events). This yields $17 \times 2 = 34$ Binomial observations.

This aggregation entails a methodological choice with two consequences. First, it discards individual-level covariates (project age and log capacity) that were retained in the static Chapter 6 analysis. The implication is that any unmodelled heterogeneity in these covariates is absorbed into the year-level intercept and slopes; results should therefore be interpreted as marginal effects averaged over the project mix observed within each (year, technology) cell. Second, the aggregation makes the three specifications nested in their likelihood structure: M1 imposes constant $\bint$ across all 17 years, M2 imposes blockwise constancy, and M3 allows $\bint$ to vary in continuous time but only along the trajectory generated by the GAS recursion. The Watanabe--Akaike Information Criterion (WAIC) and the Pareto-smoothed importance sampling Leave-One-Out (PSIS-LOO) estimator therefore yield comparable predictive-density evaluations \citep{vehtari2017practical, watanabe2010asymptotic}.

All models are implemented in PyMC \citep{salvatier2016probabilistic} with Hamiltonian Monte Carlo (No-U-Turn Sampler, \citealt{hoffman2014no}). For each model we run two chains with 2{,}000 warm-up iterations and 1{,}500 post-warm-up draws (3{,}000 effective draws after thinning), with target acceptance probability 0.98. Convergence is assessed via the $\hat{R}$ statistic \citep{vehtari2021rank} and effective sample size diagnostics. The GAS recursion is implemented via the \texttt{pytensor.scan} primitive, which allows the random-variable graph to include the recursive trajectory in a manner compatible with HMC gradient computation.

% =========================================================================
\section{Results}
\label{sec:ch7-results}

This section presents posterior inference for the three specifications. All models are fit on the year-by-technology aggregate (17 years $\times$ 2 technologies = 34 Binomial observations) on the v7 hydrogen project sample. Total events: 27 cancellations among blue projects and 16 among PEM electrolyser projects over 2010--2026.

\subsection{M1: Static baseline (aggregate)}
\label{sec:ch7-m1}

The aggregate static specification yields a posterior median $\bint = -1.37$ with 95\% credible interval $[-2.20, -0.49]$, an estimate close to the project-level Bayesian Cox PH result of $\bint = -1.43$ obtained in Chapter 6 (\textit{Spoor B-2}). The aggregate inference therefore reproduces the qualitative structural finding of a strongly negative carbon-conditional dependence, providing direct evidence that the year-by-technology aggregation does not materially distort the static result. The marginal effects of blue technology and EUA price are also recovered: $\bblue = 1.86$ [95\% CrI $1.2, 2.5$] (consistent with elevated cancellation hazard for blue at average EUA) and $\beua = 1.45$ [95\% CrI $0.76, 2.2$] (positive baseline coefficient for the EUA effect, before interaction).

\subsection{M2: Parameter-driven TVP (4-block baseline)}
\label{sec:ch7-m2}

The 4-block parameter-driven TVP specification produces the following posterior medians and 95\% credible intervals for $\bint^{b}$ across the four economic regimes:

\begin{table}[h!]
\centering
\caption{Posterior of $\bint^{b}$ across the four blocks of the parameter-driven TVP specification (M2 with $B=4$).}
\begin{tabular}{lcc}
\toprule
Block & Period & Posterior median [95\% CrI] \\
\midrule
$b_0$ & 2010--2019 (pre-energy-crisis) & $-1.50$ $[-2.40, -0.61]$ \\
$b_1$ & 2020--2022 (pandemic and early energy crisis) & $-1.73$ $[-2.80, -0.72]$ \\
$b_2$ & 2023--2024 (peak cancellations) & $-0.70$ $[-1.90, +0.41]$ \\
$b_3$ & 2025--2026 (post-peak normalisation) & $-2.06$ $[-3.80, -0.53]$ \\
\bottomrule
\end{tabular}
\label{tab:ch7-blocks}
\end{table}

Three of four blocks exhibit a credible interval that excludes zero, with point estimates clustered around $-1.5$ to $-2.1$. The second block ($b_2$, 2023--2024) is the exception: its 95\% CrI $[-1.90, +0.41]$ contains zero, and the posterior median is approximately half the magnitude of the surrounding blocks. This is the period containing 79\% of all observed cancellation events (34 of 43), so the wider posterior is not driven by data sparsity but reflects something else (Section \ref{sec:ch7-discussion} returns to this).

\subsection{M3: Observation-driven TVP (GAS, baseline $d=1/2$)}
\label{sec:ch7-m3}

The GAS specification with square-root information scaling yields posterior summaries:
\begin{align*}
\omega &= -1.61 \quad [-2.50, -0.69], \\
\phi &= 0.78 \quad [0.56, 0.94], \\
\alpha_{\mathrm{gas}} &= 0.045 \quad [0.003, 0.13], \\
\bint(0) &= -0.40 \quad [-2.10, +1.50].
\end{align*}
The long-run mean $\omega = -1.61$ is consistent with the static estimate (M1) and with the cluster of block-specific estimates from M2. The persistence parameter $\phi = 0.78$ indicates substantial autoregressive structure: in the absence of score adjustments, $\bint(t)$ would mean-revert toward $\omega$ at a rate of approximately $1 - \phi \approx 22\%$ per year. The score-response parameter $\alpha_{\mathrm{gas}} = 0.045$ is small, with the lower bound of its 95\% credible interval at 0.003. Sampling diagnostics are clean: zero divergent transitions, all $\hat R = 1.00$, all effective sample sizes above 1000.

The forward-iterated trajectory $\bint(t)$ converges to $\omega \approx -1.6$ within approximately five years and remains tightly distributed around this value: at the peak cancellation year 2023, the posterior median is $-1.57$ with 95\% CrI $[-2.59, -0.67]$, and at 2024 it is $-1.41$ with 95\% CrI $[-2.42, -0.49]$. The peak cancellation period exhibits no anomaly under the GAS specification.

\subsection{Comparison of the three specifications}

The three specifications yield qualitatively consistent inference for the carbon-conditional dependence:
\begin{itemize}
    \item All find a structurally negative $\bint$ around $-1.4$ to $-1.6$;
    \item All find that the credible interval excludes zero in the global or long-run estimate;
    \item The block specifications (M2) reveal a wider local CrI in the 2023--2024 sub-period that contains zero; the GAS specification (M3) does not exhibit this.
\end{itemize}
This combined evidence suggests that any departure from time-invariance in the carbon-conditional mechanism is small and locally confined, rather than reflecting a substantive regime shift. Section \ref{sec:ch7-discussion} returns to the methodological interpretation of this divergence between M2 and M3.

% =========================================================================
\section{Robustness Checks}
\label{sec:ch7-robust}

\subsection{Sensitivity to block boundaries}
\label{sec:ch7-blocks-robust}

We refit the parameter-driven TVP specification with $B \in \{3, 4, 5\}$ blocks. The 3-block aggregation combines pre-crisis and pandemic periods (boundary at 2022); the 5-block aggregation further splits the 2010--2019 period at 2014. Across all three specifications, the same qualitative pattern emerges: blocks centred around 2023--2024 produce wider credible intervals than blocks centred on other periods. Specifically, in the 3-block specification, $b_1$ (2023--2024) yields CrI $[-1.70, +0.50]$ (median $-0.57$); in the 4-block, $b_2$ yields $[-1.90, +0.41]$ (median $-0.70$); in the 5-block, $b_3$ yields $[-1.80, +0.39]$ (median $-0.73$). The block-boundary choice does not eliminate this localised feature, which suggests it is a robust artefact of the 2023--2024 sub-sample rather than a sensitivity to aggregation.

\subsection{Sensitivity to GAS scaling}
\label{sec:ch7-gas-robust}

The GAS specification is refit with scaling exponents $d \in \{0, 1/2, 1\}$, corresponding to no scaling, square-root information scaling (baseline), and full information scaling respectively. All three variants yield identical posterior summaries for the GAS hyperparameters ($\omega = -1.61$, $\phi = 0.78$, $\alpha_{\mathrm{gas}} = 0.045$). This is a direct consequence of the small $\alpha_{\mathrm{gas}}$: when the score-driven component contributes negligibly to the update, the scaling of the score becomes immaterial. The trajectory of $\bint(t)$ is essentially identical across the three specifications.

This robustness has an important interpretation: it directly tests the GAS framework's distinguishing feature (observation-driven updating) against the data, and finds the data does not support meaningful observation-driven dynamics. The GAS specification effectively reduces to a stationary first-order autoregression on $\bint(t)$ with mean $\omega$ and persistence $\phi$.

\subsection{Formal model comparison via LOO and WAIC}
\label{sec:ch7-icomp}

We use Pareto-smoothed importance sampling Leave-One-Out (PSIS-LOO) and Watanabe--Akaike Information Criterion (WAIC) to formally rank the seven model variants on out-of-sample predictive density of the blue-technology observation series. LOO comparison results are reported in Table \ref{tab:ch7-loo}.

\begin{table}[h!]
\centering
\caption{Pareto-smoothed leave-one-out comparison across the seven specifications (using only blue-technology observation likelihood; reported on \texttt{elpd\_loo} scale where higher is better predictive density).}
\begin{tabular}{lccc}
\toprule
Rank & Model & elpd\_loo $\pm$ se & $\Delta$elpd vs. best (dse) \\
\midrule
1 & 5-block & $-31.74 \pm 8.34$ & $0.00$ \\
2 & 3-block & $-32.13 \pm 8.52$ & $-0.39$ ($0.90$) \\
3 & 4-block & $-32.75 \pm 8.81$ & $-1.01$ ($0.82$) \\
4 & Static (M1) & $-37.06 \pm 9.97$ & $-5.32$ ($2.80$) \\
5 & GAS $d = 0$ & $-38.10 \pm 10.55$ & $-6.36$ ($4.20$) \\
6 & GAS $d = 1/2$ & $-38.10 \pm 10.55$ & $-6.36$ ($4.20$) \\
7 & GAS $d = 1$ & $-38.10 \pm 10.55$ & $-6.36$ ($4.20$) \\
\bottomrule
\end{tabular}
\label{tab:ch7-loo}
\end{table}

\textbf{Three observations}. First, the three block specifications are statistically indistinguishable from each other (all within 1 elpd unit, with dse $\leq 1$). The block boundary choice is irrelevant in formal model comparison terms. Second, the block specifications modestly outperform the static baseline ($\Delta$elpd $\approx 5$, dse $\approx 2.8$, approximately 1.9 standard errors of difference). This is suggestive but does not reach conventional thresholds of decisive evidence. Third, all three GAS variants produce \textit{identical} elpd\_loo, confirming the redundancy of the scaling parameter when $\alpha_{\mathrm{gas}}$ is near zero, and place the GAS specification approximately 6.4 elpd units below the best block specification with dse 4.2.

\textbf{Reliability caveat.} All comparisons exhibit Pareto $k > 0.7$ for some observations, indicated by the \texttt{warning: True} flag returned by the LOO computation. With 17 yearly observations the PSIS-LOO approximation is not entirely reliable, and the elpd differences should be interpreted as exploratory rather than decisive.

% =========================================================================
\section{Discussion}
\label{sec:ch7-discussion}

\subsection{What the three specifications agree on}

The static, parameter-driven block, and observation-driven GAS specifications all locate the carbon-conditional coefficient at approximately $-1.5$. The static aggregate estimate ($\bint = -1.37$) and the GAS long-run mean ($\omega = -1.61$) overlap substantially in their credible intervals. Within the parameter-driven specifications, three of four blocks (4-block specification) yield medians between $-1.5$ and $-2.1$, also overlapping. The carbon-conditional mechanism documented in Chapter 6 is therefore robust to the choice of TVP specification: it is not an artefact of the static functional form.

\subsection{Why the 2023--2024 sub-sample appears anomalous}

The parameter-driven block specifications consistently produce a wider credible interval for the 2023--2024 sub-period that includes zero, across all three block-boundary configurations. Yet the raw within-block evidence shows the expected carbon-conditional pattern: when EUA prices declined from $\bar{p}_{2023} = $\euro$82.60$ to $\bar{p}_{2024} = $\euro$65.34$ (a 21\% drop), the implied hazard ratio between blue and PEM cancellation rates rose from $1.87$ to $4.70$. Direct estimation of $\bint$ from this two-year shift yields $\bint \approx -1.28$, in line with the static and GAS estimates.

The wider posterior credible interval in the parameter-driven block specification is therefore not evidence of a regime change in the mechanism, but rather reflects the limited information available within the two-year block. The Gaussian random walk prior on $\bint^{b}$ does not pool information across blocks beyond the moderate smoothing imposed by $\sigma_{\mathrm{int}}$; with only two observations per technology in the block, the resulting posterior is necessarily diffuse. This is a known limitation of parameter-driven specifications with sparse data \citep[][Chapter 11]{durbin2012time}.

\subsection{Why GAS resolves the apparent anomaly}

The GAS specification provides a sharper estimate for $\bint(t)$ in 2023--2024 (posterior median $-1.57$ in 2023, $-1.41$ in 2024, with 95\% CrI excluding zero in both years) because it imposes structured persistence via $\phi$. Given $\phi = 0.78$, the trajectory in any year reflects an average of the long-run mean $\omega = -1.61$ and the previous year's state. The recursion mechanically borrows strength from all prior years, smoothing out any local sparsity. This is the methodological advantage of observation-driven specifications emphasised by \citet{koopman2016predicting}: regularisation through persistence rather than through explicit prior smoothing.

\subsection{Why GAS fails to outperform static in LOO}

Despite GAS providing tighter local estimates, formal LOO comparison places it below both the parameter-driven block specifications and the static baseline. The reason is two-fold. First, the small posterior estimate of $\alpha_{\mathrm{gas}}$ (0.045) implies that the GAS trajectory is dominated by the AR(1) component and does not adapt to local information much beyond what the static average captures. Second, the additional hyperparameters $\omega, \phi, \alpha_{\mathrm{gas}}, \bint(0)$ are penalised by the LOO effective parameter count without commensurate predictive improvement. The observation-driven framework provides interpretational sharpness at the cost of in-sample fit relative to the more flexible block specifications.

\subsection{Methodological lessons}

Three lessons emerge. First, in sparse-data settings, the choice between parameter-driven and observation-driven TVP specifications can lead to substantively different posterior trajectories, even when their long-run inferences coincide. Second, formal model comparison via LOO/WAIC discriminates poorly when sample size is small (17 yearly observations in our case), with Pareto-$k$ diagnostics flagging unreliability; in such settings the analyst must rely on substantive judgement about which form of regularisation is preferred. Third, the apparent regime breakdown identified by parameter-driven block specifications can be artefactual: cross-validation against an observation-driven specification with structured persistence is a useful diagnostic.

\subsection{Substantive implications}

The carbon-conditional cancellation mechanism in blue versus green hydrogen projects appears to be structurally stable across 2010--2026. The dependence of relative project survival on EUA prices does not exhibit detectable regime change despite major shocks (the 2022--2023 energy crisis, IRA introduction, ETS Phase IV transitions). For policy analysis, this stability has direct implications: forward projections of project survival under hypothetical EUA price trajectories can be made using the static estimate $\bint \approx -1.5$ without substantial loss relative to time-varying alternatives, while acknowledging that the carbon-conditional dependence is itself empirical regularity rather than a mechanistic causal model.

% =========================================================================
\section{Conclusion}
\label{sec:ch7-conclusion}

This chapter has tested the hypothesis of time-stability of the carbon-conditional hazard coefficient in hydrogen project cancellations, using three nested specifications: a static baseline, a parameter-driven random-walk TVP across economically motivated time blocks, and an observation-driven Generalised Autoregressive Score specification with score-based updating.

Three findings are summarised. First, the carbon-conditional coefficient $\bint$ is structurally negative, with posterior median estimates ranging from $-1.37$ (static) to $-1.61$ (GAS long-run mean) depending on specification, and credible intervals consistently excluding zero. Second, parameter-driven block specifications reveal a sub-period in 2023--2024 with wider credible intervals; comparison with the observation-driven GAS specification, and direct inspection of the within-block raw data, suggests this is a localised identification limitation rather than a substantive regime change. Third, formal LOO comparison places block specifications modestly above static and above GAS, but the differences are within standard errors of standard errors of difference, and the Pareto-$k$ diagnostic flags unreliable LOO estimates given the small sample.

Methodologically, this chapter demonstrates the value of comparing parameter-driven and observation-driven TVP specifications side by side in survival-style hazard analysis with sparse events, drawing on the framework of \citet{koopman2016predicting}. Substantively, it confirms the robustness of the carbon-conditional cancellation mechanism documented in Chapters 5 and 6, while showing that this mechanism does not exhibit detectable temporal variation despite the dramatic structural changes in the carbon-pricing and hydrogen-project environment during the observation window.

The chapter's main limitation is the small effective sample size. With 43 cancellation events across 17 years aggregated to 34 Binomial observations, statistical power for detecting subtle TVP signals is constrained. Future work might extend the framework with informative priors derived from the broader CCUS project sample (Chapter 8, public-data robustness analysis) or with monthly-frequency data once a sufficient observation window has accumulated.

% =========================================================================
\bibliographystyle{apalike}
\begin{thebibliography}{99}

\bibitem[Bolton and Kacperczyk(2021)]{bolton2021carbon}
Bolton, P., \& Kacperczyk, M. (2021).
\newblock Do investors care about carbon risk?
\newblock \textit{Journal of Financial Economics}, 142(2), 517--549.

\bibitem[Creal et~al.(2008)]{creal2008general}
Creal, D., Koopman, S.~J., \& Lucas, A. (2008).
\newblock A general framework for observation driven time-varying parameter models.
\newblock Tinbergen Institute Discussion Paper TI 2008-108/4.

\bibitem[Creal et~al.(2013)]{creal2013generalized}
Creal, D., Koopman, S.~J., \& Lucas, A. (2013).
\newblock Generalized autoregressive score models with applications.
\newblock \textit{Journal of Applied Econometrics}, 28(5), 777--795.

\bibitem[Durbin and Koopman(2012)]{durbin2012time}
Durbin, J., \& Koopman, S.~J. (2012).
\newblock \textit{Time Series Analysis by State Space Methods} (2nd ed.).
\newblock Oxford University Press.

\bibitem[Hoffman and Gelman(2014)]{hoffman2014no}
Hoffman, M.~D., \& Gelman, A. (2014).
\newblock The no-U-turn sampler: adaptively setting path lengths in Hamiltonian Monte Carlo.
\newblock \textit{Journal of Machine Learning Research}, 15(1), 1593--1623.

\bibitem[Koopman(2000)]{koopman2000time}
Koopman, S.~J. (2000).
\newblock Time series analysis of non-Gaussian observations based on state space models from both classical and Bayesian perspectives.
\newblock \textit{Journal of the Royal Statistical Society: Series B}, 62(1), 3--29.

\bibitem[Koopman et~al.(2016)]{koopman2016predicting}
Koopman, S.~J., Lucas, A., \& Scharth, M. (2016).
\newblock Predicting time-varying parameters with parameter-driven and observation-driven models.
\newblock \textit{Review of Economics and Statistics}, 98(1), 97--110.

\bibitem[Odenweller et~al.(2022)]{odenweller2022probabilistic}
Odenweller, A., Ueckerdt, F., Nemet, G.~F., Jensterle, M., \& Luderer, G. (2022).
\newblock Probabilistic feasibility space of scaling up green hydrogen supply.
\newblock \textit{Nature Energy}, 7(9), 854--865.

\bibitem[Salvatier et~al.(2016)]{salvatier2016probabilistic}
Salvatier, J., Wiecki, T.~V., \& Fonnesbeck, C. (2016).
\newblock Probabilistic programming in Python using PyMC3.
\newblock \textit{PeerJ Computer Science}, 2, e55.

\bibitem[Vehtari et~al.(2017)]{vehtari2017practical}
Vehtari, A., Gelman, A., \& Gabry, J. (2017).
\newblock Practical Bayesian model evaluation using leave-one-out cross-validation and WAIC.
\newblock \textit{Statistics and Computing}, 27(5), 1413--1432.

\bibitem[Vehtari et~al.(2021)]{vehtari2021rank}
Vehtari, A., Gelman, A., Simpson, D., Carpenter, B., \& B{\"u}rkner, P.~C. (2021).
\newblock Rank-normalization, folding, and localization: An improved $\hat R$ for assessing convergence of MCMC.
\newblock \textit{Bayesian Analysis}, 16(2), 667--718.

\bibitem[Watanabe(2010)]{watanabe2010asymptotic}
Watanabe, S. (2010).
\newblock Asymptotic equivalence of Bayes cross validation and widely applicable information criterion in singular learning theory.
\newblock \textit{Journal of Machine Learning Research}, 11, 3571--3594.

\end{thebibliography}

\end{document}
